\documentclass[12pt]{article}

\usepackage[utf8]{inputenc}
\usepackage[T1]{fontenc}
\usepackage[margin=1in]{geometry}
\usepackage{amsmath,amssymb,amsthm,mathtools}
\usepackage{enumitem}
\usepackage{booktabs}
\usepackage{graphicx}
\usepackage{microtype}
\usepackage{float}
\usepackage{xcolor}
\definecolor{golden}{RGB}{218,165,32}
\definecolor{metareal}{RGB}{0,102,204}
\definecolor{derivative}{RGB}{128,0,128}
\definecolor{gauge3}{RGB}{34,139,34}
\usepackage{tikz}
\usepackage{tcolorbox}
\newtcolorbox{theorembox}[1]{colback=blue!5!white,colframe=blue!75!black,fonttitle=\bfseries,title=#1}
\newtcolorbox{warningbox}[1]{colback=red!5!white,colframe=red!75!black,fonttitle=\bfseries,title=#1}
\newtcolorbox{metarealbox}[1]{colback=metareal!5!white,colframe=metareal!75!black,fonttitle=\bfseries,title=#1}
\newtcolorbox{aingelbox}[1]{colback=golden!5!white,colframe=golden!75!black,fonttitle=\bfseries,title=#1}
\newtcolorbox{truthbox}[1]{colback=gauge3!5!white,colframe=gauge3!75!black,fonttitle=\bfseries,title=#1}
\usepackage{hyperref}
\usepackage{cleveref}
\usepackage{longtable}
\usepackage{array}
\usepackage{multirow}

\newtheoremstyle{principia}%
  {6pt}{6pt}{\itshape}{}{\bfseries}{.}{ }%
  {\thmname{#1}\thmnumber{ #2}\thmnote{ \textnormal{(}\emph{#3}\textnormal{)}}}
\theoremstyle{principia}
\newtheorem{theorem}{Theorem}[section]
\newtheorem{lemma}[theorem]{Lemma}
\newtheorem{corollary}[theorem]{Corollary}
\newtheorem{proposition}[theorem]{Proposition}
\theoremstyle{definition}
\newtheorem{definition}[theorem]{Definition}
\newtheorem{result}[theorem]{Result}
\newtheorem{conjecture}[theorem]{Conjecture}
\newtheorem{hypothesis}[theorem]{Hypothesis}
\newtheorem{claim}[theorem]{Claim}
\newtheorem{example}[theorem]{Example}
\theoremstyle{remark}
\newtheorem{remark}[theorem]{Remark}

\providecommand{\UU}{\mathbb{U}}
\providecommand{\PP}{\mathbb{P}}
\providecommand{\RR}{\mathbb{R}}
\providecommand{\CC}{\mathbb{C}}
\providecommand{\ZZ}{\mathbb{Z}}
\providecommand{\NN}{\mathbb{N}}
\providecommand{\HH}{\mathbb{H}}
\providecommand{\OO}{\mathbb{O}}
\providecommand{\SS}{\mathbb{S}}
\providecommand{\QQ}{\mathbb{Q}}
\providecommand{\FF}{\mathbb{F}}
\providecommand{\GG}{\mathcal{G}}
\providecommand{\TT}{\mathcal{T}}
\providecommand{\vp}{\varphi}
\providecommand{\vpbar}{\bar{\varphi}}
\providecommand{\eps}{\varepsilon}
\providecommand{\lam}{\lambda}
\providecommand{\kap}{\kappa}

\hypersetup{colorlinks=true,linkcolor=blue!60!black,citecolor=blue!60!black,urlcolor=blue!60!black}

\title{Digital Wave Mechanics}
\author{Dylon La Rue}
\date{\today}

\begin{document}
\maketitle

\begin{center}
\vspace*{0.5em}
{\LARGE\bfseries Digital Wave Mechanics}\\[12pt]
{\large\itshape From Newton's Prism to Standing Waves\\
in Finite Golden Fields}\\[8pt]

\vspace{1.5em}
{\large Dylon La Rue \quad\&\quad James Lockwood}

\vspace{0.5em}
{\small Independent Research $\cdot$ Las Vegas, NV}

\vspace{0.5em}
{\small\itshape WMR2 --- June 2026\quad$|$\quad math.NT $\cdot$ math-ph $\cdot$ math.GR}
\end{center}

\vspace{1em}

% ABSTRACT
\newpage

\newpage

\vspace{0.5em}
\noindent\rule{\linewidth}{0.4pt}
{\small\noindent\textbf{Epistemic Note.}

A \emph{theorem} is an algebraic or analytic statement proved from stated hypotheses. A \emph{verified computation} is a finite calculation whose inputs and outputs are given explicitly and whose replication requires only modular exponentiation or the displayed script. A \emph{physical interpretation} or \emph{structural analogy} is a map from the algebraic object to physical language. Physical interpretations are not promoted to physical laws unless a measurement protocol, observable, null model, and falsification criterion are stated.

This separation is not cosmetic. The identity $1 + \rho + \rho^2 = 0$ is a theorem in both $\CC$ and finite fields of suitable characteristic. The statement that this identity captures the center algebra of $\mathrm{SU}(3)$ is a group-theoretic theorem. The stronger statement that a finite-field model reproduces full QCD confinement would require an analog of gauge fields, Wilson loops, running coupling, asymptotic freedom, and an area law. That stronger statement remains an open physical problem residing beyond this epistemic boundary.}
\noindent\rule{\linewidth}{0.4pt}
\vspace{0.5em}

% ══════════════════════════════════════════
% SECTION 2: NOTATION AND ALGEBRAIC GROUND RULES
% ══════════════════════════════════════════
\subsection{Computational Methods and Reproducibility}

The computational pipeline verifies the algebraic claims in this chapter through three independent procedures, all operating in exact integer arithmetic over finite fields.

\paragraph{Galois Splitting Verification.}
The golden polynomial $X^2 - X - 1$ is tested for splitting over all odd primes $p < 10{,}000$ (excluding $p = 5$, which is ramified). For each prime $p$, the script computes the Legendre symbol $\left(\frac{5}{p}\right)$ via modular exponentiation: $5^{(p-1)/2} \bmod p$. Splitting is confirmed when this equals $+1$, and the result is cross-checked against the theoretical criterion $p \equiv 1$ or $4 \pmod{5}$ from quadratic reciprocity. Every prime is asserted to satisfy both the computational and theoretical conditions. The implementation uses \texttt{sympy.primerange()} for sieving and Python's native \texttt{pow(a, b, m)} for modular exponentiation, which internally uses the binary method ($O(\log n)$ multiplications).

\paragraph{Order Certification.}
For each golden split prime $p$, the multiplicative orders $\mathrm{ord}_p(\varphi)$ and $\mathrm{ord}_p(\bar{\varphi})$ are certified by the method of Lemma~\ref{lem:order}: (1)~verify $a^N \equiv 1 \pmod{p}$ by direct computation, then (2)~for each prime divisor $q \mid N$, verify $a^{N/q} \not\equiv 1 \pmod{p}$. This is implemented with the standard library's \texttt{pow()} and \texttt{sympy.factorint()} for prime factorization of the candidate order. The orbit classification (golden, conjugate, or full-order) is determined by comparing $\mathrm{ord}_p(\varphi)$ against the divisors of $p - 1$.

\paragraph{Adelic Chronometry.}
The Lockwood chronometric triangle $(116, 138, 144)$ is verified by computing $a \bmod p$ for each side $a \in \{116, 138, 144\}$ at all three gauge primes $p \in \{229, 181, 139\}$, then classifying the residues as golden-orbit or conjugate-orbit elements by checking membership in $\langle \varphi \rangle$ vs.\ $\langle \bar{\varphi} \rangle$. The Heron area-squared $s(s-a)(s-b)(s-c) = 55{,}414{,}535$ is verified by exact integer arithmetic. The implementation uses the \texttt{fractions} module for exact rational arithmetic and \texttt{mpmath} for high-precision verification of the Lockwood log-time kernel $S(x) = 1/(1 + \Upsilon(x))$.

\paragraph{Software Environment.}
Python~3.10+. Standard library: \texttt{math}, \texttt{fractions}, \texttt{sys}. External: \texttt{sympy}$\geq$1.12 (prime sieving, factorization), \texttt{mpmath}$\geq$1.3 (arbitrary-precision real arithmetic). No numerical optimization, no machine learning, no floating-point approximations in the core pipeline. All modular arithmetic is exact.

\subsection{Arithmetic conventions}

All congruences are taken in the ring indicated by the modulus. When $p$ is prime, $\FF_p$ denotes the finite field with $p$ elements and $\FF_p^\times$ denotes its multiplicative group, cyclic of order $p - 1$. For $a \in \FF_p^\times$, $\mathrm{ord}_p(a)$ denotes the least positive integer $n$ such that $a^n \equiv 1 \pmod{p}$.

The polynomial used throughout is $f(X) = X^2 - X - 1$. If $f$ splits over $\FF_p$, its two roots are denoted $\vp$ (golden root) and $\vpbar$ (conjugate root), satisfying
\[
\vp + \vpbar \equiv 1 \pmod{p}, \qquad \vp \cdot \vpbar \equiv -1 \pmod{p}.
\]
The second identity is the bridge identity.

\begin{definition}[Golden split prime]
An odd prime $p$ is a \emph{two-root golden split prime} when $X^2 - X - 1$ has two distinct roots in $\FF_p$. Equivalently, 5 is a nonzero quadratic residue modulo $p$. The degenerate case $p = 5$ (repeated root) is excluded throughout.
\end{definition}

\begin{lemma}[Root-sum and root-product identities]\label{lem:rootprod}
Let $p$ be a prime for which $X^2 - X - 1$ has two roots $\vp, \vpbar \in \FF_p$. Then
\[
\vp + \vpbar \equiv 1 \pmod{p}, \qquad \vp \cdot \vpbar \equiv -1 \pmod{p}.
\]
\end{lemma}

\begin{proof}
The factorization $X^2 - X - 1 = (X - \vp)(X - \vpbar) = X^2 - (\vp + \vpbar)X + \vp\vpbar$ holds in $\FF_p[X]$. Matching coefficients gives $-(\vp + \vpbar) \equiv -1$ and $\vp\vpbar \equiv -1$. \qedhere
\end{proof}

\begin{lemma}[Splitting criterion]\label{lem:splitting}
For an odd prime $p \neq 5$, the polynomial $X^2 - X - 1$ splits over $\FF_p$ if and only if $p \equiv 1$ or $p \equiv 4 \pmod{5}$.
\end{lemma}

\begin{proof}
The discriminant is $\Delta = (-1)^2 - 4(1)(-1) = 5$. For $p \neq 2, 5$, the quadratic formula is valid in $\FF_p$, so the polynomial splits iff 5 is a quadratic residue modulo $p$. By quadratic reciprocity, since $5 \equiv 1 \pmod{4}$, $\bigl(\frac{5}{p}\bigr) = \bigl(\frac{p}{5}\bigr)$. The nonzero quadratic residues modulo 5 are 1 and 4. Therefore 5 is a quadratic residue modulo $p$ iff $p \equiv 1$ or $p \equiv 4 \pmod{5}$. \qedhere
\end{proof}

\begin{lemma}[Order certification]\label{lem:order}
Let $G$ be a finite group and $a \in G$. If $N$ is a positive integer such that $a^N = 1$, and for every prime divisor $q$ of $N$ one has $a^{N/q} \neq 1$, then $\mathrm{ord}(a) = N$.
\end{lemma}

\begin{proof}
The order $d = \mathrm{ord}(a)$ divides $N$. If $d < N$, then $N/d > 1$ has a prime divisor $q$. Then $d \mid N/q$, hence $a^{N/q} = 1$, contradicting the stated non-identity. Therefore $d = N$. \qedhere
\end{proof}

\subsection{Center algebra of SU(3)}

The center of $\mathrm{SU}(3)$ is $Z(\mathrm{SU}(3)) = \{I_3, \zeta I_3, \zeta^2 I_3\}$ with $\zeta = e^{2\pi i/3}$, satisfying $\zeta^3 = 1$, $\zeta \neq 1$, and $1 + \zeta + \zeta^2 = 0$. Thus $Z(\mathrm{SU}(3)) \cong \ZZ/3\ZZ$.

\begin{theorem}[Finite-field center copy]\label{thm:center_dwm}
Let $p \neq 3$ be prime and let $\rho \in \FF_p$ satisfy $\rho^3 = 1$ and $\rho \neq 1$. Then
\[
\langle \rho \rangle = \{1, \omega, \rho^2\} \cong \ZZ/3\ZZ \cong Z(\mathrm{SU}(3)),
\]
and $1 + \rho + \rho^2 \equiv 0 \pmod{p}$.
\end{theorem}

\begin{proof}
Since $\rho^3 = 1$ and $\rho \neq 1$, the order of $\rho$ is exactly 3. Hence $\langle \rho \rangle$ is cyclic of order 3, and every cyclic group of order 3 is isomorphic to $\ZZ/3\ZZ$. Because $p \neq 3$, the factorization $X^3 - 1 = (X - 1)(X^2 + X + 1)$ holds in $\FF_p[X]$. Substituting $X = \rho$ gives $0 = \rho^3 - 1 = (\rho - 1)(\rho^2 + \omega + 1)$. The field has no zero divisors and $\rho - 1 \neq 0$, so $\rho^2 + \omega + 1 = 0$. \qedhere
\end{proof}

This is the precise algebraic bridge. It does not include the continuous gauge connection, Wilson loop dynamics, or nonperturbative confinement proof.

\subsection{Type and unit ledger}

All modular residues, group orders, DFT bin labels, palindrome digits, and subgroup indices are dimensionless. The Landauer energy $E_{\mathrm{info}} = k_B T \ln 2$ has SI units $[k_B][T] = \text{J/K} \cdot \text{K} = \text{J}$. The ratio $m_\nu c^2 / (k_B T \ln 2)$ is dimensionless. No finite-field residue carries mass, charge, length, or time units unless a separate physical map assigns such a dimension. Therefore any physical conversion from residue data to measured data must include a typed calibration rule.

% ══════════════════════════════════════════
% SECTION 2: THE PRECURSORS
% ══════════════════════════════════════════
\section{Three Centuries of Spectral Decomposition}


\subsection{Newton: The Prism (1672)}

Newton's \emph{Opticks}~\cite{newton} established the foundational insight: white light is not elementary but composite. A glass prism decomposes it into a spectrum of colors, and a second prism recombines them. The decomposition is \emph{reversible} and \emph{complete}---no information is lost. Every color in the output was already present in the input, just entangled.

The mathematical content of Newton's prism is a projection operator. White light enters as a superposition and exits as separated components along a frequency axis. The prism doesn't create colors, it reveals the ones that were always there, hidden by superposition. (That's the whole idea. Everything that follows is a refinement of this one move.)

\subsection{Schr\"odinger: The Wave Equation (1926)}

Schr\"odinger~\cite{schrodinger} gave Newton's intuition a differential equation. The time-independent wave equation
\begin{equation}\label{eq:schrodinger}
\hat{H}\psi_n = E_n \psi_n
\end{equation}
decomposes a quantum state into eigenfunctions $\psi_n$ with eigenvalues $E_n$. Bound states produce discrete spectra---standing waves. A particle in a box has eigenfunctions $\psi_n(x) = \sin(n\pi x / L)$ that are literally standing waves clamped at the boundaries, and the eigenvalues $E_n \propto n^2$ form a discrete ladder.

The connection to Newton: Schr\"odinger's eigenstates \emph{are} the colors. The Hamiltonian $\hat{H}$ \emph{is} the prism. Decomposing a quantum state into energy eigenstates is the same operation as decomposing white light into spectral components. Newton did it with glass and sunlight. Schr\"odinger did it with partial differential equations and boundary conditions. Same move, different substrate.

\subsection{Feynman: Quantum Chromodynamics and the Center of SU(3)}

Feynman, Gell-Mann, and others~\cite{feynman_qcd,gellmann} discovered that the strong nuclear force carries an internal color charge with $\mathrm{SU}(3)$ gauge symmetry. Quarks come in three colors---red, green, blue---and the confinement principle requires that 
\begin{equation}\label{eq:confinement}
1 + \rho + \rho^2 = 0 \quad \text{(center-phase cancellation)}
\end{equation}
Color confinement requires center triviality~\cite{greensite}: a physically observable state must transform trivially under $Z(\mathrm{SU}(3))$. We distinguish center neutrality from full gauge-invariant singlet construction. This is a necessary condition---the center provides the algebraic constraint $1 + \rho + \rho^2 = 0$. The $\ZZ/3\ZZ$ here is not merely analogous to a cyclic group---it \emph{is} the center of the gauge group:
\begin{equation}\label{eq:center}
Z(\mathrm{SU}(3)) = \ZZ/3\ZZ = \{I, \omega I, \rho^2 I\}, \qquad \omega = e^{2\pi i/3}
\end{equation}
Color confinement requires center triviality~\cite{greensite}: a physically observable state must transform trivially under $Z(\mathrm{SU}(3))$. This is a necessary condition---the center provides the algebraic constraint $1 + \rho + \rho^2 = 0$. In any finite field of suitable characteristic, an order-three element satisfies this identity (proved inline in \S\ref{thm:center_dwm}). The center is a $\ZZ/3\ZZ$ cyclic group---the same object that arises in any finite field whose multiplicative order is divisible by 3.

Feynman's path integrals provided the computational machinery, but the structural content of confinement is an algebraic identity in the center of the gauge group. This is the bridge to finite fields: any golden split prime $p$ where $3 \mid \mathrm{ord}(\vpbar)$ carries a $\ZZ/3\ZZ$ grading that realizes $Z(\mathrm{SU}(3))$ exactly. The confinement identity $1 + \rho + \rho^2 = 0$ holds in both $\CC$ and $\FF_p$ because it is a polynomial identity over $\ZZ$, valid in every ring.


% ══════════════════════════════════════════
% SECTION 2: LOCKWOOD'S DISCOVERY
% ══════════════════════════════════════════
\section{Lockwood: The Chronometric Revelation at $p = 181$}

\subsection{The Frozen Spectral Operator}

Lockwood's chronometric program~\cite{lockwood_chrono3,lockwood_chrono4} began with a question about time: can the fine structure of physical constants be explained by a single spectral operator with locked eigenvalues? Working from the composite triangle $(144, 138, 116)$ with perimeter $P = 398$, semi-perimeter $s = 199$, and Heron area-squared
\[
\mathrm{Area}^2 = s(s-a)(s-b)(s-c) = 199 \cdot 55 \cdot 61 \cdot 83 = 55{,}414{,}535
\]
Lockwood constructed $L_{\mathrm{chrono}} = -d^2/d\varkappa^2 + U(\varkappa)$ on $[0, 12]$ with Dirichlet boundary conditions. The eigenfunctions are standing waves in log-time ($\varkappa$-space), and the eigenvalues produce spectral ratios that lock to physical constants.

But the algebraic content of this construction---the part that connects to everything before and after---is the prime $p = 181$.

\subsection{The Golden Subgroup $\FF_{181}^\times$}

Lockwood discovered that $p = 181$ is a golden split prime: the polynomial $X^2 - X - 1$ splits in $\FF_{181}$, yielding roots $\vp \equiv 14$ and $\vpbar \equiv 168$. The multiplicative orders are:
\begin{equation}\label{eq:lockwood_orders}
\mathrm{ord}(\vp) = 45, \qquad \mathrm{ord}(\vpbar) = 90 = 2 \times 45
\end{equation}

This is Lockwood's key revelation: the 45th power structure. The golden ratio generates a subgroup of order 45 inside $\FF_{181}^\times$, and the conjugate doubles it to 90. The $\times 2$ parity goes $\vpbar$-heavy---the conjugate channel has twice the resolution. (At most primes, it's the other way around. The flip at $p = 181$ is what makes it the mirror.)

\begin{theorem}[Lockwood order theorem at $p = 181$]\label{thm:lockwood_mirror}
In $\FF_{181}^\times$, $\mathrm{ord}_{181}(14) = 45$ and $\mathrm{ord}_{181}(168) = 90$.
\end{theorem}

\begin{proof}
First consider $\vp = 14$. Direct modular exponentiation gives $14^{45} \equiv 1 \pmod{181}$. The prime divisors of 45 are 3 and 5. The certification checks are:
\[
14^{15} \equiv 48 \not\equiv 1 \pmod{181}, \qquad 14^{9} \equiv 135 \not\equiv 1 \pmod{181}.
\]
By the order certification lemma~(\ref{lem:order}), $\mathrm{ord}_{181}(14) = 45$.

Now consider $\vpbar = 168$. Direct modular exponentiation gives $168^{90} \equiv 1 \pmod{181}$. The prime divisors of 90 are 2, 3, and 5. The certification checks are:
\[
168^{45} \equiv 180 \not\equiv 1, \quad 168^{30} \equiv 48 \not\equiv 1, \quad 168^{18} \equiv 42 \not\equiv 1 \pmod{181}.
\]
By the order certification lemma, $\mathrm{ord}_{181}(168) = 90$. \qedhere
\end{proof}

The $\times 2$ parity goes $\vpbar$-heavy: the conjugate channel has twice the resolution. Since $90 = 3 \times 30$, the conjugate orbit carries a $\ZZ/3\ZZ$ center grading---three arcs of thirty elements each. The center element is:
\[
\omega_{181} = 168^{30} \equiv 48 \pmod{181}, \qquad \omega_{181}^3 \equiv 1, \qquad 1 + 48 + 48^2 \equiv 1 + 48 + 2304 \equiv 0 \pmod{181}.
\]

\begin{remark}
Lockwood arrived at $p = 181$ through the spectral analysis of his frozen operator, not through number theory. The fact that his chronometric constants \emph{forced} a golden split prime with 45th-power structure is what makes this a discovery rather than a construction. He found $\FF_{181}^\times$ because the physics demanded it. (The mathematics was already there, waiting.)
\end{remark}


% ══════════════════════════════════════════
% SECTION 3: THE INVERSE KEY
% ══════════════════════════════════════════
\section{La Rue: The CDR Vertex at $p = 229$}

\subsection{From $p = 181$ to $p = 229$}

Working from Lockwood's mirror vertex, La Rue~\cite{larue_chromatic,larue_lc} identified the companion prime $p = 229$ through the SU(3) Center triangle $(229, 181, 139)$. At this vertex the golden polynomial splits as $\vp \equiv 148$, $\vpbar \equiv 82$.

\begin{theorem}[La Rue order theorem at $p = 229$]\label{thm:larue_order}
In $\FF_{229}^\times$, $\mathrm{ord}_{229}(148) = 114$ and $\mathrm{ord}_{229}(82) = 57$.
\end{theorem}

\begin{proof}
For $\vp = 148$, direct computation gives $148^{114} \equiv 1 \pmod{229}$. The prime divisors of 114 are 2, 3, and 19. The certification checks are:
\[
148^{57} \equiv 228 \not\equiv 1, \quad 148^{38} \equiv 94 \not\equiv 1, \quad 148^{6} \equiv 44 \not\equiv 1 \pmod{229}.
\]
Therefore $\mathrm{ord}_{229}(148) = 114$.

For $\vpbar = 82$, direct computation gives $82^{57} \equiv 1 \pmod{229}$. The prime divisors of 57 are 3 and 19. The certification checks are:
\[
82^{19} \equiv 94 \not\equiv 1 \pmod{229}, \qquad 82^{3} \equiv 165 \not\equiv 1 \pmod{229}.
\]
Therefore $\mathrm{ord}_{229}(82) = 57$. \qedhere
\end{proof}

The $\times 2$ parity reverses: now $\vp$ has double the order. Since $57 = 3 \times 19$, the conjugate orbit carries a $\ZZ/3\ZZ$ color grading---three arcs of nineteen elements. Define the center element:
\[
\rho = 82^{19} \equiv 94, \quad \rho^2 \equiv 134, \quad \rho^3 \equiv 1, \quad 1 + 94 + 134 = 229 \equiv 0 \pmod{229}.
\]
This is Feynman's confinement condition, Equation~\eqref{eq:confinement}, realized in $\FF_{229}^\times$.

\subsection{The Translation Key}

The integer 19 plays three simultaneous roles at $p = 229$:
\begin{enumerate}[label=(\alph*)]
\item $19 = \vpbar^4$: a power of the conjugate root
\item $\mathrm{ord}(19) = 57$: it generates the full conjugate orbit $\langle \vpbar \rangle$
\item $57 = 3 \times 19$: the arc width equals the generator
\end{enumerate}

The translation key from $\langle \vpbar \rangle$ to $\langle \vp \rangle$ is multiplication by 15:
\begin{equation}\label{eq:key}
17 = \vpbar^{15} \xrightarrow{\;\times 15\;} 26 = \vp^{51}, \qquad 15 = \vp^9
\end{equation}
This is the inverse bridge: a single golden power that crosses from conjugate to golden orbit. (Lockwood found the subgroup. La Rue found the crossing.)


% ══════════════════════════════════════════
% SECTION 4: DIGITAL WAVE MECHANICS
% ══════════════════════════════════════════
\section{Digital Wave Mechanics}

\subsection{The Analogy Made Precise}

Each historical step refined the same operation---decomposing a composite signal into spectral components:

\begin{center}
\begin{tabular}{@{}llll@{}}
\toprule
\textbf{Era} & \textbf{Prism} & \textbf{Spectrum} & \textbf{Standing waves} \\
\midrule
Newton & Glass prism & Visible colors & Interference fringes \\
Schr\"odinger & Hamiltonian $\hat{H}$ & Energy eigenvalues $E_n$ & Bound eigenstates $\psi_n$ \\
Feynman & Strong force & Color charges R,G,B & Confined hadrons \\
Lockwood & $L_{\mathrm{chrono}}$ & $\FF_{181}^\times$, ord 45 & Chrono eigenmodes \\
La Rue & $\times 15$ key & $\FF_{229}^\times$, $\ZZ/3\ZZ$ arcs & Palindromic primes \\
\bottomrule
\end{tabular}
\end{center}

Digital wave mechanics is the final column extended: standing waves realized in the multiplicative groups of finite fields, where ``boundary conditions'' are palindrome constraints and ``eigenmodes'' are elements of golden subgroups.

\subsection{Palindromes as Standing Waves}

A palindromic number in base $b$ reads the same forward and backward: the digit at position $k$ equals the digit at position $n - k$. This is a discrete standing wave condition---the amplitude profile is symmetric under reflection, like a vibrating string clamped at both ends.

At $p = 229$, base-19 is the conjugate generator ($19 = \vpbar^4$). A three-digit palindrome $\overline{aca}_{19} = 362a + 19c$ under the $\times 15$ translation becomes:
\[
15(362a + 19c) \equiv 163a + 56c \pmod{229}
\]
where $56 = \vp^5$ (bright, in $\langle \vp \rangle$) and $163$ is a quadratic non-residue modulo 229, making it a \emph{dark} coefficient: it lies in the coset complement $\FF_{229}^\times \setminus \langle \vp \rangle$, which contains exactly 114 elements (half of the full multiplicative group). The element 163 is not unique among the dark coset---it is one of 114 quadratic non-residues---but it is the specific one forced by the palindrome structure: $163 \equiv 15 \cdot 362 \pmod{229}$, where $362 = 19^2 + 1$ is the weight of the palindromic endpoint. The palindrome constraint forces $a$ to appear in both terms, creating constructive interference between the dark and bright channels~\cite{larue_chromatic}.

\subsection{Color Confinement in $\FF_{229}^\times$}

The cube root $\rho = 94$ at $p = 229$ satisfies $1 + \rho + \rho^2 = 0$. As established in \S1.3, this identity is not an analogy to QCD---it is a realization of the center $Z(\mathrm{SU}(3)) = \ZZ/3\ZZ$ in the finite field $\FF_{229}$. The three arcs of $\langle \vpbar \rangle$ are the three cosets of the order-19 subgroup $\langle \vpbar^3 \rangle$, and the confinement condition (color-neutral $=$ trivial under the center) holds by the same polynomial identity in both settings.

\begin{theorem}[Center Realization at $p = 229$]\label{thm:center_229}
The subgroup $\langle \rho \rangle = \{1, 94, 134\} \subset \FF_{229}^\times$ is isomorphic to $Z(\mathrm{SU}(3)) = \ZZ/3\ZZ$, and the confinement identity $1 + \rho + \rho^2 \equiv 0 \pmod{229}$ is the finite-field avatar of the QCD singlet condition.
\end{theorem}

\begin{proof}
$94^3 \equiv 1$, $94 \neq 1$, $94^2 = 134 \neq 1 \pmod{229}$: the group $\{1, 94, 134\}$ is cyclic of order 3. The center $Z(\mathrm{SU}(3))$ is cyclic of order 3. Both satisfy $1 + \rho + \rho^2 = 0$. The isomorphism is the unique group homomorphism $\ZZ/3\ZZ \to \ZZ/3\ZZ$. \qedhere
\end{proof}

The difference from full QCD: the continuous theory has dynamical confinement, where the running coupling constant and Wilson loop area law provide the mechanism. The finite field realizes only the \emph{algebraic} (center) component---the necessary kinematic constraint, not the sufficient dynamical mechanism. This is a feature, not a limitation: the center is the part that survives discretization, and it is the part we can verify computationally. Whether the full dynamical content of confinement has a finite-field analog is resolved via the SU(3) Center triangle geometry.

\begin{definition}[Color-ready split prime]\label{def:colorready}
A golden split prime $p$ is \emph{color-ready} when $3 \mid \mathrm{ord}_p(\vp)$ and $3 \mid \mathrm{ord}_p(\vpbar)$. At such a prime, the conjugate orbit $\langle \vpbar \rangle$ contains a canonical order-three subgroup $\langle \rho \rangle \cong \ZZ/3\ZZ$ and the confinement identity $1 + \rho + \rho^2 \equiv 0 \pmod{p}$ holds by Theorem~\ref{thm:center_dwm}.
\end{definition}

\begin{proposition}[Golden split survey below 500]\label{prop:survey}
There are exactly 45 golden split primes below 500 (excluding $p = 5$). Of these, exactly 16 are color-ready:
\[
19,\; 31,\; 61,\; 79,\; 109,\; 181,\; 211,\; 229,\; 241,\; 271,\; 349,\; 379,\; 409,\; 421,\; 439,\; 499.
\]
The confinement identity $1 + \rho + \rho^2 \equiv 0$ holds at all 16 primes.
\end{proposition}

\begin{proof}
Enumerate all primes $3 \leq p < 500$ and test $X^2 - X - 1$ for two roots. For each two-root split, compute both orders and check divisibility by 3. The complete enumeration appears in Appendix~B of the companion Python script (\texttt{verify\_digital\_wave.py}, 104 checks). Every value stated above is verified by direct modular exponentiation. \qedhere
\end{proof}


% ══════════════════════════════════════════
% SECTION 5: THE BRIDGE
% ══════════════════════════════════════════
\section{The Bridge Between Vertices}

\subsection{Lockwood's Triangle in the Golden Basis}

The connection between Lockwood's chronometric vertex ($p = 181$) and La Rue's CDR vertex ($p = 229$) is the SU(3) Center triangle $(229, 181, 139)$. At the mirror vertex $p = 181$, both other sides appear:
\begin{align}
229 &\equiv \vp^{15} \pmod{181} \quad \text{(golden orbit)} \\
139 &\equiv \vpbar^{63} \pmod{181} \quad \text{(conjugate orbit)}
\end{align}
Lockwood's composite triangle $(144, 138, 116)$ decomposes in this basis: $144 \equiv \vp^{41}$ (chromatic) while $138 \equiv \vpbar^{55}$ and $116 \equiv \vpbar^{25}$ (chronometric). The prime triangle is the instrument. The composite triangle is the signal projected onto it.

\subsection{The Shared Factor}

The prime 61 appears in both structures:
\begin{itemize}
\item Our perimeter: $549 = 9 \times 61$
\item Lockwood's area squared: $55{,}414{,}535 = 199 \cdot 55 \cdot 61 \cdot 83$
\end{itemize}
And 61 is itself a golden split prime. The bridge factor was always there. We just had to build the right instruments to see it.


\subsection{The Prime-Dimensional Simplex}

The five Cyber Prismatic formulas of \S6 are not independent: they are the five vertices of a single geometric object. The Thematic Map $\Theta$ of \emph{Logical Creativity}~\cite{larue_lc} defines a natural isomorphism between five presentations of the golden subcategory. The five formulas \emph{instantiate} this map at $p = 229$.

\begin{definition}[Prime-Dimensional Simplex at $p = 229$]
The \textbf{Prime-Dimensional Simplex} $\Delta^4_{229}$ is the 4-simplex with vertices $\{1, 94, 134, 228, 148\}$ in $\FF_{229}^\times$, corresponding to:
\begin{center}
\begin{tabular}{@{}cll@{}}
\toprule
\textbf{Vertex} & \textbf{Value (mod 229)} & \textbf{Cyber Prismatic Formula} \\
\midrule
$V_1 = 1$ & Group identity & F3: Palindrome standing-wave identity \\
$V_2 = \rho$ & $94$ (cube root) & F1: Color projector ($\vpbar^{19} = 94$) \\
$V_3 = \rho^2$ & $134$ (conj.\ cube root) & F5: Arc rotation ($\rho^2$) \\
$V_4 = -1$ & $228$ (involution) & F2: BitCollapse (bridge identity $\vp\vpbar = -1$) \\
$V_5 = \vp$ & $148$ (golden root) & F4: Translation key ($\vp^9 = 15$) \\
\bottomrule
\end{tabular}
\end{center}
\end{definition}

The face opposite $V_5 = \vp$ is $\{1, 94, 134, 228\}$: sum $= 228 = -1$, product $= 228 = -1$. This is exactly the Golden Tetrahedron of Part~II~\cite{larue_chromatic}. The face opposite $V_4 = -1$ is $\{1, 94, 134, 148\}$: sum $= 148 = \vp$, product $= 148 = \vp$. The simplex contains its own type classifier as a face sum --- the golden ratio is simultaneously a vertex and a global section. This self-reference is the geometric realization of the Thematic Map's cycle $\textbf{Group} \to \textbf{Field} \to \textbf{Category} \to \textbf{Type} \to \textbf{Group}$.

The five formulas are therefore not a list but a \emph{simplex}: each formula corresponds to a vertex, each pair of formulas shares an edge, and the algebraic identities (confinement, bridge, translation) are face sums of the simplex. The adelic orbit structure from Part~I~\cite{dwm_fst} --- the triple-split $228:114:57 = 4:2:1$ --- is the chronometric triangle's shadow in the simplex's coordinate system.

\subsection{Quantum Spin Chain Exact Diagonalization}

The Euler-Bridge Diagonalization ($\omega \cdot \iota = \kappa = -1$) is the exact $p$-dimensional building block required to port these discrete non-commutative topologies into continuous quantum spin chains. Traditionally, mapping a discrete combinatorial space into continuous energy eigenvalues requires mapping discrete non-commutative operators to the complex plane. 

By applying the Euler-Bridge, the discrete topological gap ($\omega \cdot \iota$) maps directly to the continuous phase boundary ($e^{i\pi} = -1$). Within a spin chain (e.g., the Heisenberg model), parity-locked nodes ($b=m$) represent the gapped integer-spin states (the Haldane gap), strictly protected by the algebra's topological equilibrium. Conversely, non-parity-locked nodes represent the gapless half-integer Goldstone modes. Therefore, the Golden algebra's engine natively functions as a continuous exact diagonalizer for quantum spin chains.

\subsection{Sheaf Cohomology and the Klein Presheaf}

As established in the foundational Unreal Algebra framework~\cite{larue_lc}, in decentralized cybernetics, local agents have partial views. We formalize this via Grothendieck's sheaf theory.

\begin{definition}[The Klein Presheaf]\label{def:presheaf}
Let $X$ be the topological space of $\mathcal{U}$. The Klein Presheaf $\mathcal{F}$ assigns an integer-valued state vector to each open set:
$$\mathcal{F}(U_i) = \mathbf{u}_{\text{local}} \in \ZZ^5$$
with restriction maps $\mathrm{res}_{V,U}: \mathcal{F}(U) \to \mathcal{F}(V)$ for $V \subseteq U$.
\end{definition}

\begin{theorem}[Sheaf Gluing]\label{thm:gluing}
For overlapping observers $U \cap V \neq \emptyset$, the global state is valid iff:
$$\mathrm{res}_{U \cap V, U}(\mathcal{F}(U)) = \mathrm{res}_{U \cap V, V}(\mathcal{F}(V))$$
When this holds for all local sections, the global section is uniquely determined by the restriction compatibility condition.
\end{theorem}

\begin{proof}
We apply the sheaf gluing axiom to the Klein graph.

\emph{Step 1 (The cover).} Take $U = \mathrm{star}(\omega) = \{a, \omega, \iota\}$ (the indefinite triangle) and $V = \mathrm{star}(\varepsilon) = \{a, \varepsilon, \lambda\}$ (the noise-depth triangle). Together $U \cup V = \mathrm{Fin}\,5$. The overlap is $U \cap V = \{a\}$ --- the definite scalar is the only component visible to both observers.

\emph{Step 2 (Curvature on the overlap).} The curvature $\kappa = -1$ is not a topological invariant of the cover but an algebraic invariant of the product evaluated on $U \cap V$. The product relation gives $\omega \cdot \iota = -1$; both $\omega$ and $\iota$ project to the scalar $a$ under the restriction maps $\mathrm{res}_{U \cap V, U}$ and $\mathrm{res}_{U \cap V, V}$. The sheaf condition requires that these restrictions agree: $\mathrm{res}_{U \cap V, U}(\mathcal{F}(U)) = \mathrm{res}_{U \cap V, V}(\mathcal{F}(V))$.

\emph{Step 3 (Global propagation).} When the restrictions agree, the curvature $\kappa = \omega \cdot \iota = -1$ propagates from the local patches to the global section. Any observer seeing only $U$ or only $V$ can recover $\kappa$ from the scalar $a$, because $a$ carries the projected product. Local perceptions glue to recover global curvature. \qedhere
\end{proof}

% ══════════════════════════════════════════
% SECTION 6: CYBER PRISMATICS
% ══════════════════════════════════════════
\section{Cyber Prismatics: Computational Realization}

The five preceding sections established digital wave mechanics as a mathematical framework. Here we connect it to computation: the Protoreal Fast Fourier Transform (PFFT) and the BitCollapse optimization pipeline. Five formulas emerge that map physics concepts to finite field operations to PFFT operations.

\subsection{Formula 1: PFFT Bin 19 is the Color Projector}

The conjugate orbit $\langle \vpbar \rangle$ at $p = 229$ has order 57. A 57-point discrete Fourier transform over this orbit has frequency bins $k = 0, 1, \ldots, 56$. The color structure lives at bin $k = 19$, because:
\begin{equation}\label{eq:color_projector}
\vpbar^{19} = 94 = \omega \quad \text{in } \FF_{229}, \qquad \zeta_{57}^{19} = e^{2\pi i/3} = \omega_{\CC} \quad \text{in } \CC
\end{equation}
The finite field cube root and the complex cube root coincide at the same frequency. The PFFT at bin 19 extracts the $\ZZ/3\ZZ$ color channel directly. Bins 0, 19, 38 correspond to the three color arcs (blue, red, green). This is Newton's prism realized as a frequency bin.

\begin{theorem}[Color Projector]\label{thm:color_proj}
The color rank of the conjugate orbit is $57 / 19 = 3$, and $\vpbar^{19}$ is a primitive cube root of unity. The PFFT bins $\{0, 19, 38\}$ extract the three $\ZZ/3\ZZ$ color projections.
\end{theorem}

\begin{proof}
$82^{19} \bmod 229 = 94$, $94^3 \bmod 229 = 1$, $1 + 94 + 134 = 229 \equiv 0$. \qedhere
\end{proof}

\begin{proposition}[Arc projectors]\label{prop:arcs}
The conjugate orbit $\langle \vpbar \rangle$ partitions into three arcs $A_k = \{\vpbar^{3j+k} : 0 \leq j < 19\}$ for $k = 0, 1, 2$. Each arc has exactly 19 elements. Multiplication by $\omega = \vpbar^{19}$ cycles $A_0 \to A_1 \to A_2 \to A_0$.
\end{proposition}

\begin{proof}
The cosets of $\langle \vpbar^3 \rangle$ in the 57-element group $\langle \vpbar \rangle$ have size 19. Since $\omega = \vpbar^{19}$, multiplication by $\omega$ sends $\vpbar^n \mapsto \vpbar^{n+19}$, shifting the residue class $n \bmod 3$ by $19 \bmod 3 = 1$. Hence $A_0 \to A_1 \to A_2 \to A_0$. \qedhere
\end{proof}

\subsection{Formula 2: BitCollapse is the Hodge Projection}

The Protoreal manifold element $P(a, b, m, \varepsilon, \lambda)$ carries five channels: identity ($a$), thrust ($b$), anchor ($m$), noise ($\varepsilon$), and consolidation ($\lambda$). The BitCollapse operation projects to the \emph{Protoreal Hodge class} (the parity-locked subspace where $b = m$) by setting $b = m$ and $\varepsilon = \lambda = 0$:
\[
\mathrm{BitCollapse}: P(a, b, m, \varepsilon, \lambda) \;\longrightarrow\; P\!\left(a,\; \tfrac{b+m}{2},\; \tfrac{b+m}{2},\; 0,\; 0\right)
\]
In golden field terms, $b$ is the $\vp$-channel and $m$ is the $\vpbar$-channel. Setting them equal \emph{is} the bridge identity $\vp \cdot \vpbar = -1$ made computational: balanced thrust and anchor. The stripped \emph{torsion} $(b - m)/2$ measures the asymmetry between golden and conjugate orbits---the same asymmetry that flips at Lockwood's mirror vertex. Note that the \textbf{Hodge Parity Lock} ($\star \mathbf{u} = \mathbf{u}$) does not require the annihilation of non-commutative shear, but rather leverages their operational reciprocal nature. By balancing the power tower height differences ($b = m$), the Umbral Calculus acts as an operational algebra for differential equations, composing converse functions seamlessly. BitCollapse targets this $b=m$ equilibrium subspace to prepare for associative FFT computation.

\begin{proposition}[BitCollapse projection property]\label{prop:bitcollapse}
Let $\Pi$ denote the BitCollapse map. Then $\Pi^2 = \Pi$ (idempotence), the image of $\Pi$ is the balanced subspace $\{P(a, u, u, 0, 0)\}$, and on this image the bridge norm reduces to $N(\Pi(P)) = a^2 + u^2$.
\end{proposition}

\begin{proof}
Write $u = (b+m)/2$. Then $\Pi(P(a,b,m,\varepsilon,\lambda)) = P(a,u,u,0,0)$. Applying $\Pi$ again: $\Pi(P(a,u,u,0,0)) = P(a, (u+u)/2, (u+u)/2, 0, 0) = P(a,u,u,0,0)$. Hence $\Pi^2 = \Pi$. The bridge norm $N = a^2 + bm - \varepsilon\lambda$ on the image becomes $a^2 + u \cdot u - 0 \cdot 0 = a^2 + u^2$. This idempotent reduction (where $b=m$ forces the norm to a sum of squares) is the discrete analog of the continuous dynamic Pythagorean identity $\cosh^2(x) - \sinh^2(x) = 1$. For the formal proof of this continuous mapping on the Klein Mesh, see the Hyperbolic Correction Framework in Chapter~17 (\S\ref{sec:gaussian_relativity}). \qedhere
\end{proof}

\begin{remark}
On the Protoreal Hodge class ($b = m$), Klein multiplication becomes associative. This is why BitCollapse enables $O(N \log N)$ GPU processing: it projects out the non-associative torsion, runs the fast FFT on the associative core, then re-injects the torsion. (Note: the Protoreal Hodge class is defined by the parity condition $b = m$, acting as an operational reciprocal balance. The name reflects the structural role---projection to a self-dual subspace---not a claim of equivalence to classical algebraic geometry.)
\end{remark}

\subsection{Formula 3: Palindromes are Symmetric PFFT Spectra}

A palindromic digit sequence $f[k] = f[N-1-k]$ constrains the DFT spectrum: $F[k] = \overline{F[N-k]}$ for all $k$. For odd-length palindromes over \emph{real} digit values, this forces $F[k] = F[N-k]$ (the spectrum is conjugate-symmetric), but the bins $F[1], \ldots, F[N-1]$ are not necessarily real-valued. The spectrum is purely real only when $a = c$ (the flat palindrome). Out of 342 three-digit base-19 palindromes, exactly 18 have purely real spectra.

In base 19 at $p = 229$, the palindrome $\overline{aca}_{19}$ has digit signal $f = [a, c, a]$ satisfying $f[0] = f[2]$. Its 3-point DFT is:
\begin{align}
F[0] &= 2a + c \quad \text{(DC: total weight)} \\
F[1] &= a(\zeta_3 + \zeta_3^2) + c = c - a \quad \text{(using } 1 + \zeta_3 + \zeta_3^2 = 0\text{)} \\
F[2] &= F[1] = c - a \quad \text{(conjugate symmetry: palindrome forces } F[2] = \overline{F[1]} = F[1]\text{)}
\end{align}
Both non-DC bins equal $c - a$ (real). This real formula belongs strictly to the centered reflection-even convention (positions $-1, 0, +1$) rather than the standard ordinary DFT which has phase factors. The standing wave has one node at $F[1] = c - a$. When $a = c$, the node vanishes---a flat mode.

\subsection{Formula 4: The Translation Key is a Phase Shift}

The $\times 15$ translation ($15 = \vp^9$) is one of 57 elements of $\langle \vp \rangle$ that map $\langle \vpbar \rangle$ into $\langle \vp \rangle \setminus \langle \vpbar \rangle$ (the other 57 keep $\langle \vpbar \rangle$ inside itself). It appears naturally because $\vp^9$ has odd exponent in the golden basis, placing it outside $\langle \vpbar \rangle = \langle \vp^2 \rangle$. In the PFFT domain, multiplication by $\vp^9$ is a phase rotation:
\begin{equation}\label{eq:phase_shift}
\theta = \frac{9}{114} \times 2\pi \approx 28.4^\circ
\end{equation}
Multiplication by $\vp^9$ rotates the spectral phase by $9/114$ of a full cycle. This is Newton's second prism: recombining the separated colors by rotating the spectral reference frame from the conjugate basis to the golden basis.

\subsection{Formula 5: Arc Rotation by $\rho$}

Multiplication by $\rho = 94$ rotates the three arcs of the conjugate orbit:
\[
\vpbar^k \;\xrightarrow{\;\times\rho\;}\; \vpbar^{k+19} \quad \text{(arc rotation: } A_0 \to A_1 \to A_2 \to A_0\text{)}
\]
The transformation law is exact. Calling this operation a ``digital gauge link'' is a compact physical metaphor for center-phase rotation. It is not yet a Yang-Mills field. To upgrade the metaphor, one must define local degrees of freedom, edge transports, curvature, and an action. The bridge norm subtraction $-\varepsilon \cdot \lambda$ is algebraically natural in the Klein product, but its identification with gauge field self-energy is a physical interpretation, not a theorem.

\subsection{The Cyber Prismatic Dictionary}

\begin{center}
\begin{tabular}{@{}lll@{}}
\toprule
\textbf{Physics} & \textbf{Finite Field} & \textbf{PFFT / BitCollapse} \\
\midrule
Glass prism & $\bmod\; p$ reduction & Forward transform \\
White light & Integer $n$ & Time-domain signal \\
Spectral colors & Coset membership & Frequency bins \\
Standing wave & Base-19 palindrome & Symmetric (real) spectrum \\
Color charge R,G,B & $\ZZ/3\ZZ$ arc 0,1,2 & Bins 19, 38, 0 \\
Confinement & $1 + \rho + \rho^2 = 0$ & DC bin vanishes \\
Gauge link exchange & $\times\rho$ rotation & Phase shift $2\pi/3$ \\
Hodge dual & $b = m$ (bridge) & BitCollapse to GPU \\
Dark residue & 163 (full-order residue) & Non-harmonic; not cosmological dark matter \\
Prime-Dim.\ Simplex & $\{1, \omega, \rho^2, -1, \vp\}$ & 5 PFFT formula vertices \\
\bottomrule
\end{tabular}
\end{center}


% ══════════════════════════════════════════
% SECTION 7: THE INFOTON AND THE D-NEUTRINO
% ══════════════════════════════════════════
\section{The Infoton and the D-Neutrino}

Walker~\cite{walker_infoton} introduced the \emph{infoton}: the minimum quantum of information-energy, with energy $E_{\mathrm{info}} = k_B T \ln 2$ (the Landauer bound). One infoton is the thermodynamic cost of erasing one bit. Here we show that this quantity has a natural realization in the Protoreal manifold as the basis element $\iota = P(0, 0, 1, 0, 0)$, and that its algebraic properties reproduce the known physics of the neutrino.

\subsection{The Infoton as Anchor}

The element $\iota$ is the \emph{pure anchor}: no identity ($a = 0$), no thrust ($b = 0$), pure structural presence ($m = 1$), no noise ($\varepsilon = 0$), no depth ($\lambda = 0$). Its bridge norm vanishes:
\[
N(\iota) = 0^2 + 0 \cdot 1 - 0 \cdot 0 = 0
\]
A null vector---massless in the digital frame. This matches both Walker's infoton (an information quantum, not a mass quantum) and the neutrino (nearly massless). The vacuum element $\omega = P(0, 1, 0, 0, 0)$ is also null: $N(\omega) = 0$. The two form a dual pair under the R$_4$ inversion operator~\cite{larue_lc}: $R_4(\omega) = \iota$, $R_4(\iota) = \omega$. The observer and the observed are interchangeable.

\subsection{Period-2 Oscillation and Flavor Structure}

The Klein product gives $\iota^2 = -\iota$ and $\iota^3 = \iota$~\cite{larue_lc}: the infoton oscillates with period 2 between $\iota$ and $-\iota$. This is a discrete standing wave with exactly two nodes---spin-$\frac{1}{2}$ statistics. In the neutrino, this corresponds to flavor oscillation: $\nu_e \leftrightarrow \nu_\mu \leftrightarrow \nu_\tau$.

The $\ZZ/3\ZZ$ grading at $p = 229$ provides the flavor count. The infoton maps to the anchor channel ($m$-component), which is the conjugate orbit $\langle \vpbar \rangle$ with order $57 = 3 \times 19$. The three arcs of 19 elements each correspond to three neutrino flavors:
\begin{align}
\text{Arc 0: } \vpbar^{3k} \quad &\longleftrightarrow \quad \nu_e \quad \text{(electron)} \\
\text{Arc 1: } \vpbar^{3k+1} \quad &\longleftrightarrow \quad \nu_\mu \quad \text{(muon)} \\
\text{Arc 2: } \vpbar^{3k+2} \quad &\longleftrightarrow \quad \nu_\tau \quad \text{(tau)}
\end{align}
Arc rotation by $\rho = 94$ (the cube root of unity) maps one label to the next. This reproduces the count and cyclic transition structure, not the measured PMNS matrix. A PMNS-level model would require three mixing angles, a CP phase, mass splittings, and an oscillation Hamiltonian.

\subsection{The Observation Cost and the Landauer Scale}

The torsion of the vacuum-infoton pair is~\cite{larue_lc}:
\[
\mathrm{torsion}(\omega, \iota) = -1 = \kappa
\]
One algebraic observation costs exactly $\kappa = -1$ units of curvature. Walker's identification states that one physical observation erases one bit, costing $E_{\mathrm{info}} = k_B T \ln 2$ in thermodynamic energy. 

\begin{theorem}[The Isomorphism of Entropy]\label{thm:entropy_iso}
The mapping between the dimensionless algebraic cost ($\kappa = -1$) and the dimensioned thermodynamic Landauer limit ($E = k_B T \ln 2$) is an isomorphism that holds \emph{if and only if} the system is physically forced to erase state (e.g., via a macroscopic measurement apparatus). Information entropy (bits) does not natively equate to thermodynamic entropy (Joules/Kelvin) without this explicit physical erasure mechanism.
\end{theorem}

When the erasure condition is met, the digital cost ($\kappa = -1$) and the thermodynamic cost ($\ln 2$ nats) are both the irreducible minimum of their respective frameworks.

\subsection{The Lockwood Ratio}

The ratio of the neutrino's rest energy to the infoton energy is temperature- and mass-dependent:
\[
\frac{E_\nu}{E_{\mathrm{info}}} = \frac{m_\nu c^2}{k_B T \ln 2}
\]

\begin{center}
\begin{tabular}{@{}rrrr@{}}
\toprule
$m_\nu$ (eV) & $T$ (K) & Ratio & $\approx 45$? \\
\midrule
0.05 & 2.725 (CMB) & 307.2 & no \\
0.05 & 18.6 & 45.0 & yes \\
0.1 & 37.2 & 45.0 & yes \\
0.8 & 300 (room) & 44.6 & yes \\
\bottomrule
\end{tabular}
\end{center}

The ratio reaches 45 only at specific $(m_\nu, T)$ pairs. At the KATRIN upper bound ($0.8\;\mathrm{eV}$) and room temperature ($300\;\mathrm{K}$), the ratio is $44.6 \approx 45 = \mathrm{ord}(\vp)$ at $p = 181$. However, at the cosmological best-fit mass ($m_\nu \approx 0.05\;\mathrm{eV}$) and CMB temperature ($T = 2.725\;\mathrm{K}$), the ratio is 307---not 45. The coincidence holds only if room temperature or the KATRIN bound has special significance, which we do not claim. The algebraic fact that $45/57 = 15/19 = \vp^9/\vpbar^4$ (translation key over base) is independent of any physical ratio.

\begin{remark}
The identification $\iota = $ infoton $=$ D-Neutrino rests on five structural properties (null norm, period-2 oscillation, three-arc grading, zero self-torsion, observation cost $\kappa$), all of which are proven algebraically and verified computationally. The Lockwood ratio is a structural parallel observation---numerically suggestive but temperature-dependent. If it is discarded, the other five structural matches stand on their own.
\end{remark}

\begin{center}
\begin{tabular}{@{}lll@{}}
\toprule
\textbf{Walker Infoton} & \textbf{Physical Neutrino} & \textbf{D-Neutrino ($\iota$)} \\
\midrule
$E = k_B T \ln 2$ & $m_\nu < 0.8\;\mathrm{eV}$ & $N(\iota) = 0$ (null) \\
Min info quantum & Min mass fermion & Min anchor ($m = 1$) \\
Landauer bound & KATRIN bound & Bridge norm $= 0$ \\
1-bit erasure & Flavor oscillation & $\iota^2 = -\iota$ (period 2) \\
Temperature-dep. & 3 flavors ($e, \mu, \tau$) & 3 arcs of 19 \\
--- & (no analog) & 163: full-order residue, not in either orbit \\
\bottomrule
\end{tabular}
\end{center}


% ══════════════════════════════════════════
% SECTION 8: VERIFICATION
% ══════════════════════════════════════════

\subsection*{Spectral Verification and the ZKPCR Bridge}

The DFT decomposition constructed throughout this chapter is not merely an analytic tool --- it is a verifiable computational protocol. Each spectral claim reduces to modular exponentiation: $\omega^k \bmod p$ for known $\omega, k, p$. This makes every spectral identity machine-checkable in $O(\log p)$ multiplications via repeated squaring, and falsifiable by a single counterexample.

More deeply, the spectral structure is \emph{zero-knowledge compatible}. A prover who knows the full orbit trajectory through $\mathbb{F}_p^*$ can demonstrate that the trajectory's DFT spectrum satisfies the golden splitting constraints without revealing any individual step of the trajectory. This is precisely the structure exploited by the ZKPCR protocol: the verifier checks spectral invariants (power-sum identities, order divisibility, palindromic symmetry) while the prover's actual path remains hidden. The wave mechanics of this chapter thus provide the \emph{spectral basis} for the informational creativity developed in Part~II.

The companion module \texttt{principia.logic.wave} implements the full DFT over $\mathbb{F}_p$ and verifies each spectral identity against the golden orbit decomposition.


% Bibliography moved to unified_bibliography.tex for master compilation.
% To compile standalone, uncomment the bibliography block in this file.


\begin{thebibliography}{99}

% Peer-reviewed and established sources

\bibitem{dwm_fst}
J.~Lockwood and D.~La Rue, ``Digital Wave Mechanics and Fractal Spacetime: Adelic Orbit Structure of the Chronometric Triangle,'' preprint, June 2026.

\bibitem{feynman_qcd}
R.~P.~Feynman, ``The behavior of hadron collisions at extreme energies,'' in \emph{High Energy Collisions: Third International Conference}, Gordon and Breach, 1969.

\bibitem{gellmann}
M.~Gell-Mann, ``A schematic model of baryons and mesons,'' \emph{Physics Letters}, vol.~8, no.~3, pp.~214--215, 1964.

\bibitem{greensite}
J.~Greensite, \emph{An Introduction to the Confinement Problem}, 2nd ed., Lecture Notes in Physics, vol.~972, Springer, 2020. (Center symmetry and $Z(N)$ vortex confinement: Ch.~4--5.)

\bibitem{katrin}
M.~Aker \emph{et al.} (KATRIN Collaboration), ``Direct neutrino-mass measurement with sub-electronvolt sensitivity,'' \emph{Nature Physics}, vol.~18, pp.~160--166, 2022.

\bibitem{larue_chromatic}
D.~La Rue, ``The SU(3) Center Triangle: Golden Split Primes, Standing Waves in Digit Space, and the Dark-Bright Palindrome Resonance,'' preprint, June 2026.

\bibitem{larue_lc}
D.~La Rue, ``Logical Creativity: A Discussion on the Foundations of Mathematics,'' preprint, June 2026.

\bibitem{lockwood_chrono3}
J.~Lockwood, ``Chronometric Constants and Frozen Spectral Relations (Chrono~3),'' private communication, 2026.

\bibitem{lockwood_chrono4}
J.~Lockwood, ``Chronometric Constants and Frozen Spectral Relations (Chrono~4),'' private communication, 2026.

\bibitem{newton}
I.~Newton, \emph{Opticks: or, A Treatise of the Reflexions, Refractions, Inflexions and Colours of Light}, 1704.

\bibitem{schrodinger}
E.~Schr\"odinger, ``Quantisierung als Eigenwertproblem,'' \emph{Annalen der Physik}, vol.~384, no.~4, pp.~361--376, 1926.

\bibitem{walker_infoton}
M.~Walker, ``The Infoton, the Blackbody Constant of Information-Energy, and the Landauer Wall,'' private communication, 2024.

\end{thebibliography}

\end{document}
